\section{一の冪根で残る中心と表現の構造}
\label{sec:notes-root-unity-v38}

\subsection{中心局所化を使う理由}

端点の全変形族 $S_H=B[w;\sigma,x^{-1}[\,-,H]]$ では、
一般の $H$ に対して $x$ は正規元でない。
しかし対角パラメータ $t$ が位数 $\ell$ の一の冪根なら、
$\Xi=x^\ell$ は $Z(B)$ に入り、$\sigma$ が固定し、混合Ore導分も零にする。
従って $\Xi$ は全代数の中心元となる。
これにより
\[
 S_H[\Xi^{-1}]=B[x^{-1}][J],\qquad J=xw+H
\]
を正当な中心局所化として使える。
この手順は、正規元がOre拡大で正規性を失う場合にも、有限位数の冪が
中心性を回復すれば適用できる。

\subsection{周期が中心の特異点と表現次元を同時に決める}

$\beta=(4\gamma)^{-1}$、$U=u^\ell$、$V=v^\ell$ とすれば、
中心の唯一の関係は
\[
 UV=\prod_{i=0}^{\ell-1}(-Q-t^{2i+1}\beta x^2).
\]
$t^2$ の位数が $\ell/\gcd(2,\ell)$ であることから、
\[
 \begin{array}{c|c|c}
 \ell&\text{中心の関係}&\text{特異中心指標の単純加群}\\ \hline
 \text{奇数}&UV=-Q^\ell-\beta^\ell\Xi^2&\text{特異指標なし}\\
 2d&UV=(Q^d+(-1)^d\beta^d\Xi)^2&d\text{ 次元が二つ}
 \end{array}
\]
となる。ここでは $\Xi$ を可逆としている。
奇数位数では全中心上でAzumaya、偶数位数では中心の滑らかな部分が正確にAzumaya部分である。
この結論は一般的なAzumaya部分と滑らかな部分の一致定理に依存しない。
$x$ の固有値を分離する有限エタール被覆上で、周期quiverの可逆辺を縮約して証明できる。

偶数位数の特異ファイバーは $M_d(\mathcal R)$ であり、$\mathcal R$ は
二頂点・各向き二本の矢を持つ平方零代数である。
その全次元 $6d^2$、半単純商の次元 $2d^2$、根基の次元 $4d^2$ は区別して残す。
PI次数は $2d$ のままであり、特異点の単純加群の次元 $d$ と取り違えない。

\subsection{別の中心開部分への移行}

混合項 $H$ の全ての方向は、関数斜体だけでなく上の中心局所化からも消える。
v39では $U=u^\ell$ を反転する別の中心開部分を用い、$\Xi=0$ への延長が
$H\bmod x$ の巡回ノルムを記憶することを決定した。
全中心 $Z(S_H)=Z(B[x^{-1}])[J]\cap S_H$ の決定には、なお
$\Xi=U=V=0$ の深い境界を扱う必要がある。
非次数付き局所化の同型から、全次数付き環や射影圏の同値を推論しない。
根基の平方零模型は、境界上の単純加群の退化や拡張を比較する際の基準模型になる。
